\documentclass{article}
\date{}

\usepackage[norsk]{babel}
\usepackage{subcaption}
\usepackage{amsmath}
\usepackage{graphicx}
\usepackage[a4paper, top=1in, bottom=1in, left=1in, right=1in]{geometry}
\usepackage{subcaption}
\usepackage{amsfonts}
\usepackage{dirtytalk}
\usepackage{enumerate}

% vertical vector = \vect{...}
\newcommand{\vect}[1]{\begin{pmatrix} #1 \end{pmatrix}}
% double underline is \uu{...}
\newcommand{\uu}[1]{\underline{\underline{#1}}}
% double bold underline is \uub{...}
\newcommand{\uub}[1]{\underline{\underline{\textbf{#1}}}}
% align counter reset
\AtBeginEnvironment{align}{\setcounter{equation}{0}}
% arrows
\newcommand{\lra}{\quad\Longrightarrow\quad}
% equivalent = \eq 
\newcommand{\eq}{\quad\Longleftrightarrow\quad}
% 3x space
\newcommand{\trip}{\ \ \ }

\setlength{\parskip}{1em}
\setlength{\parindent}{0pt}

\begin{document}

\section*{Seksjon 4.1}

\subsection*{Oppgave 1}
Løsninger til ligningsystemet:
\begin{align}
    x + 2y - z = 3 \\
    2x + 3y - 3z = -1 \\
    -x + 2y + 3z = 1 
\end{align}

(1), (2), og (3) betegner hver linje.

Trinn 1 (isoler x): \\ \\
$(2) = (2) + (3)\cdot2 \lra [2x+3y-3z=-1]\ +\ [-2x+4y+6z=2] \lra [7y+3z=1]$ \\ 
$(3) = (3) + (1) \quad\, \lra [-x+2y+3z=1]\trip + \ [x+2y-z=3] \quad\ \ \lra [4z+2z=4]$


\setcounter{equation}{0}
\begin{align}
    x + 2y - z = 3 \\
    7y + 3z = 1 \\
    4y + 2z = 4
\end{align}

Trinn 2 (isoler y): \\ \\
$(3) = 7\cdot(3)\ - 4\cdot(2) \lra [28y+14z=28]\ -\ [28y+12z=4] \lra [2z=24]$

\setcounter{equation}{0}
\begin{align}
    x + 2y -z = 3 \\
    7y + 3z = 1 \\
    2z = 24
\end{align}

Løser: \\\\
\begin{align*}
    2z = 24 \eq z = 12 \\
    7y + 3\cdot12 = 1 \eq 7y = -35 \eq y = -5 \\
    x + 2\cdot-5 -(12) = 3 \eq x-10+12 = 3 \eq x = 25
\end{align*}

\[\uu{z = 12}, \qquad \uu{y = -5}, \qquad \uu{x = 25}\]

\newpage
\subsection*{Oppgave 3}
Løsninger til ligningsystemet:

\begin{align}
    2x - 4y + 6z = -2 \\
    -3x + 2y -z = 8 \\
    x - 6y + 11z = 4
\end{align}

Trinn 1 (isoler x): \\\\
$(1) = (1) - 2\cdot(3) \lra [2x-4y+6z=-2]\ +\ [-2x+12y-22z=-8] \lra [8y-16z = -10]$

$(2) = (2) + 3\cdot(3) \lra [-3x+2y-z=8]\ \ \, +\ [3x-18y+33z=12]\quad \lra [-16y+32z = 20]$ 

$(3) = (1)$ og $(1) = (3) \qquad$ (bytter 1 og 3)

$(3) = (2)$ og $(2) = (3)$ \qquad (bytter 2 og 3)

\begin{align}
    x - 6y + 11z = 4 \\
    8y - 16z = -10 \\
    -16y + 32z = 20 
\end{align}

Trinn 2 (isoler y): \\\\
$(3) = (3) + 2\cdot(2) \lra [-16y+32z=20] + [16y-32z = -20] \lra [0 = 0]$

\begin{align}
    x-6y+11z=4 \\
    8y-16z=-10 \\
    0 = 0
\end{align}

Løser: \\\\
\begin{align}
    0 = 0 \\
    8y-16z = -10 \eq 8y = 16z-10 \eq y = 2z - \frac{5}{4} \\
    x - 12z+\frac{15}{2} + 11z = 4 \eq x = z - \frac{7}{2} 
\end{align}

\[
\uu{Z = fri variabel}, \qquad \uu{y= 2z-\frac{5}{4}}, \qquad \uu{x = z-\frac{7}{2}}
\]
\newpage
\subsection*{Oppgave 4}
Løsninger til ligningsystemet:

\begin{align}
    x-2y+3z = 1 \\
    -x+y-2z=0 \\
    -3x+5y-8z=2
\end{align}

Trinn 1 (isoler x): \\\\
$(2) = (2) + (1) \trip \, \lra [-x+y-2z=0]\quad \ +\ [x-2y+3z=1]\ \, \lra [-y+z=0]$ \\
$(3) = (3) + 3\cdot(1) \lra [-3x+5y-8z=2]\ +\ [3x-6y+9z=3] \lra [-y+z=5]$

\begin{align}
    x-2y+3z=1 \\
    -y+z = 0 \\
    -y+z=5
\end{align}

Løser: \\\\
(2) og (3) motsier hverandre, så dette ligningsystemet går ikke.

\section*{Seksjon 4.2}

\subsection*{Opppgave 1}

\begin{enumerate}[A=]
    \item \( \begin{pmatrix} 1 & -2 & 3 \\ 0 & 1 & -1 \\ 0 & 0 & 1 \end{pmatrix} \) - Ja
    \item \( \begin{pmatrix} 0 & 1 & 2 \\ 0 & 0 & 1 \\ 0 & 0 & 0 \end{pmatrix} \) - Ja
    \item \( \begin{pmatrix} 1 & 0 & 2 \\ 0 & 2 & 0 \\ 0 & 0 & 1 \end{pmatrix} \) - Nei (Pivot rad 2 er 2, ikke 1)
    \item \( \begin{pmatrix} 1 & -2 & 0 & -1 \\ 0 & 0 & 1 & 1 \\ 0 & 0 & 0 & 1 \end{pmatrix} \) - Ja
    \item \( \begin{pmatrix} 1 & 3 & -2 \\ 0 & 1 & -3 \\ 0 & 0 & 1 \\ 0 & 0 & 0 \end{pmatrix} \) - Ja
    \item \( \begin{pmatrix} 1 & -2 & 0 & -1 \\ 0 & 1 & 2 & 1 \\ 0 & 1 & 4 & -3 \end{pmatrix} \) - Nei (Pivot rad 3 er i samme kolonne som rad 2)
\end{enumerate}





\subsection*{Opppgave 2}

\begin{enumerate}[A=]
    \item \( 
    \begin{pmatrix}
        1&2&1\\-1&-3&2\\2&1&2
    \end{pmatrix}\)

    \item \(
    \begin{pmatrix}
        0&1&2\\1&1&2\\-2&1&3
    \end{pmatrix}\)

    \item \(
    \begin{pmatrix}
        1&1&-2&3\\2&1&3&0\\-1&0&-5&2
    \end{pmatrix}
    \)

    \item \(
    \begin{pmatrix}
        1&2&1&1\\-1&3&-1&2\\1&12&1&7
    \end{pmatrix}
    \)
    
\end{enumerate}




\subsection*{Opppgave 3}





\subsection*{Opppgave 4}


\end{document}

